\section{Tolman--Eichenbaum Machine: What, Where y Associative Memory}

TEM es el núcleo del modelo de la segunda parte de la clase. Hace una analogía entre la lateral entorhinal cortex (LEC) y el content/what stream, entre la medial entorhinal cortex (MEC) y el structural/location stream, y entre el hippocampus y la conjunctive associative memory. El estado estructural se actualiza de forma recurrente (recurrently) según la action, el contenido se vincula rápidamente en un entorno nuevo, y la memory sostiene después el recall y la inference.

\subsection{Panorama de la arquitectura: cada uno de los tres estados cumple su función}

Después de la portada, el diagrama de la arquitectura muestra tres capas de información: sensory observation, structural state reutilizable y conjunctive memory. El graph de la parte inferior indica que una misma estructura puede alojar distintas observation labels. El modelo no mete la observation directamente en una RNN, sino que hace que la structure dynamics y el content encoding se aprendan por separado y luego se combinen en un hippocampal-like state.

\lecturefigure{slide-50-tem-title.jpg}{Tolman--Eichenbaum Machine: unificación de la memoria espacial y la relacional}{Estado didáctico manuscrito del video oficial V050; 00:59:38}

\lecturefigure{slide-51-tem-architecture.jpg}{Arquitectura de TEM: sensory content, structural state y memory}{Estado didáctico manuscrito del video oficial V051; 01:01:08}

\begin{longtable}{@{}L{0.17\textwidth}L{0.31\textwidth}L{0.40\textwidth}@{}}
\toprule
Componente & Función computacional & Analogía neurocientífica \\
\midrule
$x_t$ / sensory code & Representación del contenido de la observation actual & LEC-like what/content stream \\
$g_t$ / structural code & Representación de posición/relación actualizada a partir de la acción y el estado anterior & MEC-like where/structure stream \\
$p_t$ / conjunctive code & Vincula el contenido actual con la estructura en un episode & hippocampal conjunctive representation \\
$M_t$ / associative memory & Escribe rápidamente los conjunctive states y los recupera a partir de un cue & hippocampal recurrent/associative memory \\
\bottomrule
\end{longtable}

\subsection{LEC y MEC: separación entre contenido y estructura}

La diapositiva de LEC extrae el sensory code a partir de la observation y responde «qué se vio»; la diapositiva de MEC actualiza el structural code según la action/transition y responde «en qué lugar de la estructura se está». Esta separación permite que el modelo conserve la transition dynamics en un entorno nuevo y, al mismo tiempo, sustituya la object identity. Si ambas estuvieran completamente entrelazadas desde el inicio, cada cambio de observation podría destruir el mapa existente.

\lecturefigure{slide-52-lec-what-stream.jpg}{LEC-like stream: codificación del sensory content}{Estado didáctico manuscrito del video oficial V052; 01:01:20}

\lecturefigure{slide-53-mec-where-stream.jpg}{MEC-like stream: actualización del reusable structural state}{Estado didáctico manuscrito del video oficial V053; 01:02:24}

La actualización estructural puede abstraerse como
$$
g_t=f_{\theta}(g_{t-1},a_t),
\qquad
x_t=e_{\phi}(o_t),
$$
donde $o_t$ es la observation, $a_t$ es la action/edge, $g_t$ es el path-integrated structural state y $x_t$ es el content encoding. $\theta$ se aprende lentamente a través de los entornos, mientras que el object binding del entorno actual puede establecerse con rapidez.

\begin{warningbox}{«Where» no es necesariamente una coordenada bidimensional}
En una tarea de tree o de transitive-order, $g_t$ representa una relational position y no tiene que corresponder a un $(x,y)$ euclidiano. Interpretar el MEC-like state como un código estructural general es justamente lo que permite a TEM generalizar de la navigation a la relational memory; sin embargo, esa generalización todavía requiere el respaldo de tareas y de datos neuronales.
\end{warningbox}

\subsection{Hippocampal binding y Hebbian-style memory}

La diapositiva de memory-update combina content y structure en un conjunctive code, que después se escribe en la associative memory. La diapositiva Hebbian ofrece la intuición de una actualización local en la que «la activación conjunta refuerza la conexión». Una forma abstracta de escribirlo es
$$
p_t=b(x_t,g_t),
\qquad
M_t=\lambda M_{t-1}+p_t p_t^{\top},
$$
donde $b$ es la binding operation, $p_t$ es la conjunctive representation, $M_t$ es la memory matrix y $\lambda$ controla cuánto se conserva la memoria antigua. El TEM real usa una estructura de inference/generative más detallada, pero esta forma simplificada muestra por qué la escritura rápida no exige volver a entrenar todos los parámetros.

\lecturefigure{slide-54-recurrent-memory-update.jpg}{Recurrent state y associative-memory update en TEM}{Estado didáctico manuscrito del video oficial V054; 01:02:58}

\lecturefigure{slide-55-hebbian-memory-equations.jpg}{Hebbian-style binding y memory equations}{Estado didáctico manuscrito del video oficial V055; 01:03:40}

\begin{knowledgebox}{Qué son los fast weights}
Los parámetros del modelo $\theta,\phi$ se entrenan lentamente a lo largo de muchos entornos, mientras que la memory matrix $M_t$ se actualiza rápidamente dentro del episode actual; por eso suele llamársele fast weights. Esto permite guardar contenido nuevo tras una sola exposure o unas pocas, mientras la dynamics estructural sigue proviniendo de los parámetros de largo plazo.
\end{knowledgebox}

\subsection{Graph inference y planning}

La diapositiva de graph-inference muestra cómo el agent usa el structural state y la recalled memory para inferir nodos que no observó directamente o para planear rutas. Lo esencial no es almacenar el graph completo como una adjacency table fija, sino hacer que el recurrent transition model genere una structure candidata y que después la memory restituya la observation identity del entorno actual.

\lecturefigure{slide-56-graph-inference-plan.jpg}{Graph inference a partir del structural state y la memory}{Estado didáctico manuscrito del video oficial V056; 01:06:02}

\begin{lstlisting}[caption={Flujo conceptual del estado estructural y la vinculación rápida}]
def tem_step(structure, action, observation, memory):
    next_structure = transition_model(structure, action)
    content = content_encoder(observation)
    conjunctive = bind(next_structure, content)
    memory.write(conjunctive)
    recalled = memory.read(next_structure)
    return next_structure, recalled
\end{lstlisting}

\teachervoice{Will describe el objetivo como «la estructura se aprende lentamente; el contenido se vincula rápidamente». Esto hace eco de la separación address/value de SDM en la primera mitad: la structure es como un address system transferible, el sensory content es como el value que se escribe, y la hippocampal memory combina ambos de forma temporal.}

\subsection{Grid-like code y adaptación rápida}

Las dos diapositivas de grid-code comparan la representation del modelo con el periodic spatial pattern que suele aparecer en los datos neuronales; la diapositiva de learning curve muestra que el agent mejora rápidamente después de pocas visitas. La conclusión correcta es que un recurrent state estructurado, junto con el entrenamiento en la tarea, puede producir una grid-like basis y sostener un map learning rápido. Eso no significa que cada hidden unit corresponda a una grid cell real, ni demuestra que la retícula hexagonal sea la única codificación viable.

\lecturefigure{slide-57-grid-code-evidence.jpg}{Grid-like representations que surgen en TEM}{Estado didáctico manuscrito del video oficial V057; 01:08:26}

\lecturefigure{slide-58-grid-code-summary.jpg}{Comparación estructural entre el modelo y los neural grid codes}{Estado didáctico manuscrito del video oficial V058; 01:09:36}

\lecturefigure{slide-59-quick-learning-result.jpg}{Curva de aprendizaje que mejora rápidamente tras visitar pocos nodos}{Estado didáctico manuscrito del video oficial V059; 01:10:48}

\begin{knowledgebox}{Lectura de la figura: grid-like no se juzga solo por «si parece hexagonal»}
Hay que examinar la periodicidad, la orientación, la escala y la fase del firing field, así como su estabilidad ante cambios del entorno, y comparar la population structure. Las manchas repetidas en un solo mapa de calor pueden deberse al suavizado, a los bordes o a los datos de entrenamiento; solo cuando varias estadísticas geométricas coinciden con la causalidad de la tarea se tiene una evidencia más sólida.
\end{knowledgebox}

\subsection{Resumen de la sección}

TEM logra un aprendizaje rápido de mapas mediante un content stream, un structural recurrent state y una fast associative memory. La representación estructural se reutiliza entre entornos, las observation nuevas se vinculan rápidamente y la memory sostiene la graph inference. El grid-like code y el few-visit learning son evidencia funcional, pero no deben presentarse como una identidad anatómica.
